\begin{prob}{2.9}{}{京大院試 H2 数学I-4}
有理数上で無理数、無理数上で有理数の値を取る連続関数$f:\mathbb{R}\to\mathbb{R}$は存在するか？
\end{prob}

$f(\mathbb{R}) = f(\mathbb{Q}) \cup f(\mathbb{R} \setminus \mathbb{Q}) \subseteq f(\mathbb{Q}) \cup \mathbb{Q}$
であるから, 像は高々可算無限集合である.
ここでもし像が異なる2点$a,b$を持てば, 連続性(中間値の定理)から像は区間$[a,b]$を含む.
しかし$[a,b]$は非可算であるため, 不適. よって像は1点のみ, すなわち定数関数となるが, これは条件を満たさない.
実際, その定数が有理数であるならば「有理数上で無理数値を取る」は成り立たず,
またその定数が無理数であるならば「無理数上で有理数値を取る」は成り立たないからである.

以上より, このような連続関数は存在しない. 
